\section{柱坐标系下的分离变量法}
\idx{柱坐标}{zhuzuobiao}
\label{sec:cylindrical}

圆形边界在直角坐标里是 $x^2+y^2=a^2$.若试 $V=X(x)Y(y)$,边界上 $X(x)Y(\sqrt{a^2-x^2})=0$ 拆不成只约束 $X$ 的条件和只约束 $Y$ 的条件.改用极坐标以后,同一条边界就是 $\rho=a$,乘积解的两个因子可以各自去满足边界.

与 $z$ 无关时,\eqref{eq:laplace-cyl}退化为
\begin{equation}
  \frac{\partial^2 V}{\partial\rho^2}
  +\frac{1}{\rho}\frac{\partial V}{\partial\rho}
  +\frac{1}{\rho^2}\frac{\partial^2 V}{\partial\phi^2}
  =0.
  \label{eq:polar-laplace}
\end{equation}
设 $V=R(\rho)\Phi(\phi)$,除以 $R\Phi$ 再乘 $\rho^2$,
\begin{equation}
  \frac{\rho}{R}\frac{\mathrm{d}}{\mathrm{d}\rho}
    \Bigl(\rho\frac{\mathrm{d}R}{\mathrm{d}\rho}\Bigr)
  =-\frac{1}{\Phi}\frac{\mathrm{d}^2\Phi}{\mathrm{d}\phi^2}
  =m^2.
\end{equation}
左右只依赖各自的自变量,所以等于同一个常数.这里写成 $m^2$,是为了让角向方程的周期解取正频率.

$\phi$ 增加 $2\pi$ 回到同一点,$\Phi(\phi+2\pi)=\Phi(\phi)$.于是
\begin{equation}
  \Phi''+m^2\Phi=0,
\end{equation}
$m=0,1,2,\ldots$,本征函数是 $1$,以及 $\cos(m\phi)$ 和 $\sin(m\phi)$.$m$ 不是任意实数,单值性把它离散化,作用和固定端把弦的波数离散化相同.

径向方程是欧拉方程
\begin{equation}
  \rho^2 R''+\rho R'-m^2 R=0.
\end{equation}
令 $\rho=e^{s}$,则 $\rho R_{\rho}=R_s$, $\rho^2 R_{\rho\rho}=R_{ss}-R_s$,代入得 $R_{ss}-m^2 R=0$.
$m\geq 1$ 时
\begin{equation}
  R=A\rho^{m}+C\rho^{-m}.
\end{equation}
$m=0$ 时 $R_{ss}=0$,
\begin{equation}
  R=C_0+D_0\ln\rho.
\end{equation}
$\rho^{m}$ 在无穷远处增长,$\rho^{-m}$ 和 $\ln\rho$ 在原点发散.区域包不包原点、包不包无穷远,决定哪一组系数必须为零.

外部区域 $\rho>a$ 的一般解是
\begin{equation}
\begin{aligned}
  V={}&C_0+D_0\ln\rho \\
  &+\sum_{m=1}^{\infty}\Bigl[
      (A_m\rho^{m}+C_m\rho^{-m})\cos(m\phi)
      +(B_m\rho^{m}+D_m\rho^{-m})\sin(m\phi)
    \Bigr].
\end{aligned}
\label{eq:polar-general}
\end{equation}
内部区域 $\rho<a$ 则丢掉 $\ln\rho$ 和 $\rho^{-m}$.

\subsection{均匀电场中的导体圆柱}

\begin{figure}[htb]
  \centering
  \begin{tikzpicture}[font=\small, >=Stealth,
    vec/.style={-{Stealth[length=2mm]}, line width=0.7pt}]
  \draw[gray!35, thin] (-2.6,0) -- (2.6,0) node[right] {$x$};
  \draw[gray!35, thin] (0,-1.9) -- (0,1.9) node[above] {$y$};
  \fill[blue!8] (0,0) circle (0.85);
  \draw[blue!70!black, thick] (0,0) circle (0.85);
  \node at (0,0) {导体};
  \foreach \y in {-1.35,-0.45,0.45,1.35} {
    \draw[vec, gray!80] (-2.45,\y) -- (-1.15,\y);
    \draw[vec, gray!80] (1.15,\y) -- (2.45,\y);
  }
  \draw[vec, red!70!black] (0.85,0) -- (1.55,0) node[midway, above] {$2E_0$};
  \draw[vec, red!70!black] (-1.55,0) -- (-0.85,0);
  \node[font=\scriptsize, gray] at (2.05,1.55) {$\mathbf{E}_0$};
  \node[font=\scriptsize] at (0,-1.55) {$r=a$};
\end{tikzpicture}

  \caption{均匀电场中的中性导体圆柱.远处场恢复为 $\mathbf{E}_0$,表面与外场对齐的两点法向场增强到 $2E_0$.}
  \label{fig:cylinder-field}
\end{figure}

半径为 $a$ 的无限长导体圆柱放在垂直于轴的均匀电场 $\mathbf{E}_0$ 里.取轴为 $z$,$x$ 轴平行于 $\mathbf{E}_0$.导体以外没有电荷,势满足\eqref{eq:polar-laplace}.导体是等势体,把这个常数选成零,
\begin{equation}
  V(a,\phi)=0.
\end{equation}
远处 $E_x\to E_0$, $E_y\to 0$,均匀场的势是 $-E_0 x=-E_0\rho\cos\phi$.若圆柱单位长度还带电 $q_0$,远场再加一项对数,
\begin{equation}
  V(\rho,\phi)\sim
  -\frac{q_0}{2\pi\varepsilon_0}\ln\rho
  -E_0\rho\cos\phi
  \qquad(\rho\to\infty).
\end{equation}
可加常数已经由导体零势定掉.

把\eqref{eq:polar-general}放上 $V(a,\phi)=0$.傅里叶系数各自为零,
\begin{equation}
  C_0+D_0\ln a=0,
  \qquad
  C_m=-a^{2m}A_m,
  \qquad
  D_m=-a^{2m}B_m.
\end{equation}
每一个增长的径向幂都配上一个衰减项,使它们在 $\rho=a$ 上抵消.

远场最高只到 $\rho^1$,所以 $m\geq 2$ 的 $A_m,B_m$ 必须是零,连带的 $C_m,D_m$ 也是零.远场没有 $\sin\phi$,故 $B_1=D_1=0$.余弦的一次项要再现 $-E_0\rho\cos\phi$,故 $A_1=-E_0$, $C_1=E_0 a^2$.对数项的系数是 $D_0=-q_0/(2\pi\varepsilon_0)$,于是 $C_0+D_0\ln\rho=D_0\ln(\rho/a)$.合起来
\begin{equation}
  V(\rho,\phi)
  =\frac{q_0}{2\pi\varepsilon_0}\ln\frac{a}{\rho}
  -E_0\Bigl(\rho-\frac{a^2}{\rho}\Bigr)\cos\phi,
  \qquad\rho\geq a.
  \label{eq:cylinder-potential}
\end{equation}

\begin{example}
中性圆柱的表面电场.
\end{example}
\begin{solution}
$q_0=0$ 时
\begin{equation}
  V=-E_0\Bigl(\rho-\frac{a^2}{\rho}\Bigr)\cos\phi.
\end{equation}
在 $\rho=a$ 上,$\rho-a^2/\rho=0$,导体保持零势.$\rho\to\infty$ 时修正项消失,$V\sim -E_0\rho\cos\phi$.径向电场
\begin{equation}
  E_{\rho}=-\frac{\partial V}{\partial\rho}
  =E_0\Bigl(1+\frac{a^2}{\rho^2}\Bigr)\cos\phi.
\end{equation}
表面上 $E_{\rho}=2E_0\cos\phi$.$\phi=0$ 和 $\phi=\pi$ 两点,法向场的大小都是背景场的两倍,电击穿首先从这里开始.$\phi=\pm\pi/2$ 时 $\cos\phi=0$,中性圆柱的势在这两条射线上处处等于导体的势.
\end{solution}

\subsection{圆内的边值}

\begin{example}
圆盘 $\rho<a$ 的边界上 $V(a,\phi)=V_0\cos\phi$,盘内无电荷,势在圆心有限.
\end{example}
\begin{solution}
内部解丢掉负幂和对数,
\begin{equation}
  V=A_0+\sum_{m=1}^{\infty}\rho^{m}(A_m\cos(m\phi)+B_m\sin(m\phi)).
\end{equation}
边界只有 $\cos\phi$,所以只留下 $m=1$ 的余弦项,$A_1 a=V_0$,
\begin{equation}
  V(\rho,\phi)=V_0\frac{\rho}{a}\cos\phi=V_0\frac{x}{a}.
\end{equation}
它是直角坐标里的线性函数,拉普拉斯算子当然是零.圆形边界把这个简单解从傅里叶级数里挑了出来.
\end{solution}

\subsection{有限圆柱}
$\rho$ 和 $z$ 都有边界时,分离常数落到贝塞尔方程上.考虑稳态温度:圆柱 $0\leq\rho<a$, $0\leq z\leq L$,侧面和底面保持零度,顶面给定轴对称分布 $f(\rho)$,
\begin{equation}
  \nabla^2 T=0,
  \qquad
  T(a,z)=0,
  \qquad
  T(\rho,0)=0,
  \qquad
  T(\rho,L)=f(\rho).
\end{equation}
$f$ 与 $\phi$ 无关,解也与 $\phi$ 无关.\eqref{eq:laplace-cyl}变成
\begin{equation}
  \frac{1}{\rho}\frac{\partial}{\partial\rho}
    \Bigl(\rho\frac{\partial T}{\partial\rho}\Bigr)
  +\frac{\partial^2 T}{\partial z^2}=0.
\end{equation}
设 $T=R(\rho)Z(z)$.为了让 $R(a)=0$ 而且 $R$ 在轴上有限,分离常数取正,
\begin{equation}
  \frac{1}{\rho R}\frac{\mathrm{d}}{\mathrm{d}\rho}
    \Bigl(\rho\frac{\mathrm{d}R}{\mathrm{d}\rho}\Bigr)
  =-k^2
  =\frac{1}{Z}\frac{\mathrm{d}^2 Z}{\mathrm{d}z^2}.
\end{equation}
径向方程是 $0$ 阶贝塞尔方程.令 $\xi=k\rho$,则
\begin{equation}
  \xi^2\frac{\mathrm{d}^2 R}{\mathrm{d}\xi^2}
  +\xi\frac{\mathrm{d}R}{\mathrm{d}\xi}
  +\xi^2 R=0,
\end{equation}
在 $\xi=0$ 有限的解是 $J_0(k\rho)$.侧面 $J_0(ka)=0$,所以
\begin{equation}
  k_m=\frac{j_{0,m}}{a},
  \qquad m=1,2,\ldots,
\end{equation}
$j_{0,m}$ 是 $J_0$ 的第 $m$ 个正零点.轴向方程 $Z''-k_m^2 Z=0$,底面 $Z(0)=0$ 选出
\begin{equation}
  Z=\sinh(k_m z).
\end{equation}
叠加
\begin{equation}
  T(\rho,z)
  =\sum_{m=1}^{\infty} A_m J_0(k_m\rho)\sinh(k_m z).
\end{equation}
顶面 $z=L$ 用贝塞尔族的正交性定系数.权重为 $\rho$ 的模是
\begin{equation}
  \int_0^{a}\rho\, J_0(k_m\rho)^2\,\mathrm{d}\rho
  =\frac{a^2}{2}J_1(k_m a)^2,
\end{equation}
因此
\begin{equation}
  A_m\sinh(k_m L)
  =\frac{2}{a^2 J_1(k_m a)^2}
    \int_0^{a} f(\rho)\, J_0(k_m\rho)\,\rho\,\mathrm{d}\rho.
  \label{eq:cylinder-coeff}
\end{equation}

\begin{example}
顶面温度是常数 $T_0$.
\end{example}
\begin{solution}
用 $\mathrm{d}[u J_1(u)]/\mathrm{d}u=u J_0(u)$,
\begin{equation}
  \int_0^{a}\rho\, J_0(k_m\rho)\,\mathrm{d}\rho
  =\frac{a}{k_m}J_1(k_m a).
\end{equation}
代入\eqref{eq:cylinder-coeff},
\begin{equation}
  A_m
  =\frac{2T_0}{k_m a\, J_1(k_m a)\,\sinh(k_m L)}.
\end{equation}
于是
\begin{equation}
  T(\rho,z)
  =\sum_{m=1}^{\infty}
    \frac{2T_0}{j_{0,m} J_1(j_{0,m})\, \sinh(k_m L)}
    \, J_0(k_m\rho)\sinh(k_m z),
\end{equation}
其中用了 $k_m a=j_{0,m}$.$z=0$ 时每一项都是零;$z=L$ 时双曲正弦之比变成 $1$,剩下的级数就是常数 $T_0$ 的贝塞尔展开.侧面 $\rho=a$ 上 $J_0(k_m a)=0$.
\end{solution}
